\documentclass[a4paper, 10pt]{article}
% Per-subject config for this repo. See vendor/template/course-config.tex.example
% for what each macro is used for.

\newcommand{\CourseCode}{2110515}
\newcommand{\CourseName}{Intro to Quantum Computing}
\newcommand{\StudentID}{6733172621}
\newcommand{\StudentName}{Patthadon Phengpinij}
\newcommand{\DefaultCollaborators}{Claude (for\,\LaTeX\,styling and grammar checking)}

\usepackage{../vendor/template/CEDT-Assignment-style}

\begin{document}
\subject
\asgntitle{3}{}{Week 3}{6733172621 Patthadon Phengpinij}{Claude (for\,\LaTeX\,styling and grammar checking)}

% ================================================================================ %
%                                    Problem 01                                    %
% ================================================================================ %
\begin{problem}
\textbf{Finding Eigenvectors by Hand}
\par This is a tedious exercise.
But I want you to go through to practice the skill of finding eigenvectors.
We can also measure the electron spin in \( \hat{y} \) direction by setting up an external magnetic field in the \( \hat{y} \) direction.
It is given that the corresponding operator is \[ \sigma_y = \begin{pmatrix}
	0 & -i \\
	i & 0
\end{pmatrix}. \] Again, we are using the \( \ket{0}/\ket{1} \) basis.
Find the eigenvalues and eigenvectors in the \( \ket{0}/\ket{1} \) basis.
\end{problem}

\begin{solution}
	To find the eigenvalues, we solve the characteristic equation \( \det\paren{\sigma_y - \lambda I} = 0 \):
	\[
		\det\begin{pmatrix}
			-\lambda & -i       \\
			i        & -\lambda
		\end{pmatrix} = \lambda^2 - (-i)(i) = \lambda^2 - 1 = 0
	\]
	which gives \( \lambda_1 = 1 \) and \( \lambda_2 = -1 \).

	\textbf{Eigenvector for \( \lambda_1 = 1 \).} Solve \( \paren{\sigma_y - I}\vec{v} = \vec{0} \):
	\[
		\begin{pmatrix}
			-1 & -i \\
			i  & -1
		\end{pmatrix}
		\begin{pmatrix}
			v_1 \\ v_2
		\end{pmatrix} = \begin{pmatrix}
			0 \\ 0
		\end{pmatrix}
	\]
	The first row gives \( -v_1 - i v_2 = 0 \), i.e. \( v_1 = -i v_2 \). Setting \( v_1 = 1 \) gives \( v_2 = i \), so the unnormalized eigenvector is \( \begin{pmatrix} 1 \\ i \end{pmatrix} \). Normalizing (norm \( = \sqrt{1^2 + |i|^2} = \sqrt{2} \)):
	\[
		\ket{\lambda_1 = 1} = \frac{1}{\sqrt{2}}\begin{pmatrix}
			1 \\ i
		\end{pmatrix} = \frac{1}{\sqrt{2}}\paren{\ket{0} + i\ket{1}}
	\]

	\textbf{Eigenvector for \( \lambda_2 = -1 \).} Solve \( \paren{\sigma_y + I}\vec{v} = \vec{0} \):
	\[
		\begin{pmatrix}
			1 & -i \\
			i & 1
		\end{pmatrix}
		\begin{pmatrix}
			v_1 \\ v_2
		\end{pmatrix} = \begin{pmatrix}
			0 \\ 0
		\end{pmatrix}
	\]
	The first row gives \( v_1 - i v_2 = 0 \), i.e. \( v_1 = i v_2 \). Setting \( v_1 = 1 \) gives \( v_2 = -i \), so the unnormalized eigenvector is \( \begin{pmatrix} 1 \\ -i \end{pmatrix} \). Normalizing:
	\[
		\ket{\lambda_2 = -1} = \frac{1}{\sqrt{2}}\begin{pmatrix}
			1 \\ -i
		\end{pmatrix} = \frac{1}{\sqrt{2}}\paren{\ket{0} - i\ket{1}}
	\]

	\textbf{Check.}
	\[
		\sigma_y \ket{\lambda_1 = 1} = \frac{1}{\sqrt{2}}\begin{pmatrix}
			0 & -i \\
			i & 0
		\end{pmatrix}\begin{pmatrix}
			1 \\ i
		\end{pmatrix} = \frac{1}{\sqrt{2}}\begin{pmatrix}
			1 \\ i
		\end{pmatrix} = \ket{\lambda_1 = 1}
	\]
	\[
		\sigma_y \ket{\lambda_2 = -1} = \frac{1}{\sqrt{2}}\begin{pmatrix}
			0 & -i \\
			i & 0
		\end{pmatrix}\begin{pmatrix}
			1 \\ -i
		\end{pmatrix} = -\frac{1}{\sqrt{2}}\begin{pmatrix}
			1 \\ -i
		\end{pmatrix} = -\ket{\lambda_2 = -1}
	\]
	both of which confirm \( \sigma_y \ket{\lambda_k} = \lambda_k \ket{\lambda_k} \), as expected.
\end{solution}
% ================================================================================ %

% ================================================================================ %
%                                    Problem 02                                    %
% ================================================================================ %
\begin{problem}
\textbf{Finding Eigenvectors Using Python}
\par Use Google Colab to check your answer. What you need is the function to find the
eigenvectors and eigenvalues. The following code first builds the matrix and then find the
eigenvectors and eigenvalues:
\begin{codingbox}
sigmay = np.array([[0, -1j], [1j, 0]])
print(sigmay)

eigenvalues, eigenvectors = np.linalg.eig(sigmay)
print(eigenvalues)
print(eigenvectors)
\end{codingbox}
\end{problem}

\begin{solution}
	Running the provided code in Google Colab prints
	\[
		\text{eigenvalues} \approx \begin{pmatrix}
			1 \\ -1
		\end{pmatrix}, \qquad
		\text{eigenvectors} \approx \begin{pmatrix}
			-0.7071i & 0.7071   \\
			0.7071   & -0.7071i
		\end{pmatrix}
	\]
	where each column of the eigenvector matrix is the eigenvector for the eigenvalue in the corresponding position.

	\begin{enumerate}
		\item  The first column, \( \begin{pmatrix} -0.7071i \\ 0.7071 \end{pmatrix} = -i \cdot \frac{1}{\sqrt{2}}\begin{pmatrix} 1 \\ i \end{pmatrix} \),
		is the hand-calculated eigenvector for \( \lambda_1 = 1 \) up to the global phase factor \( -i \),
		which is physically irrelevant since an eigenvector is only defined up to a nonzero scalar multiple.

		\item The second column, \( \begin{pmatrix} 0.7071 \\ -0.7071i \end{pmatrix} \approx \frac{1}{\sqrt{2}}\begin{pmatrix} 1 \\ -i \end{pmatrix} \),
		matches the hand-calculated eigenvector for \( \lambda_2 = -1 \) exactly.
	\end{enumerate}
\end{solution}
% ================================================================================ %

% ================================================================================ %
%                                    Problem 03                                    %
% ================================================================================ %
\begin{problem}
\textbf{Pauli Matrix Multiplication by Hand and Using Python}
\par Use hand calculation to show that \( \sigma_y^2 = I \).
Then verify with Colab using the following codes.
Note that \texttt{np.matmul} is for matrix multiplication.
\begin{codingbox}
sigmay = np.array([[0, -1j], [1j, 0]])

print(sigmay)
print(np.matmul(sigmay, sigmay))
\end{codingbox}
\end{problem}

\begin{solution}
	By hand:
	\[
		\sigma_y^2 = \begin{pmatrix}
			0 & -i \\
			i & 0
		\end{pmatrix}\begin{pmatrix}
			0 & -i \\
			i & 0
		\end{pmatrix} = \begin{pmatrix}
			(-i)(i) & 0       \\
			0       & (i)(-i)
		\end{pmatrix} = \begin{pmatrix}
			-i^2 & 0    \\
			0    & -i^2
		\end{pmatrix} = \begin{pmatrix}
			1 & 0 \\
			0 & 1
		\end{pmatrix} = I
	\]

	Verifying with the provided Python code, the output is
	\begin{align*}
		\sigma_y &= \begin{pmatrix}
			0 & -i \\
			i & 0
		\end{pmatrix} \\
		\sigma_y^2 &= \begin{pmatrix}
			1 & 0 \\
			0 & 1
		\end{pmatrix}
	\end{align*}
	
	which matches the hand-calculated result \( \sigma_y^2 = I \).
\end{solution}
% ================================================================================ %

% ================================================================================ %
%                                    Problem 04                                    %
% ================================================================================ %
\begin{problem}
\textbf{Pauli Matrix Multiplication by Hand and Using Python}
\par Repeat the previous question for \( \sigma_z^2 = I \) using both hand calculation and Colab.
\end{problem}

\begin{solution}
	By hand, with \( \sigma_z = \begin{pmatrix} 1 & 0 \\ 0 & -1 \end{pmatrix} \):
	\[
		\sigma_z^2 = \begin{pmatrix}
			1 & 0  \\
			0 & -1
		\end{pmatrix}\begin{pmatrix}
			1 & 0  \\
			0 & -1
		\end{pmatrix} = \begin{pmatrix}
			1^2 & 0      \\
			0   & (-1)^2
		\end{pmatrix} = \begin{pmatrix}
			1 & 0 \\
			0 & 1
		\end{pmatrix} = I
	\]

	Verifying with Colab:
\begin{codingbox}
sigmaz = np.array([[1, 0], [0, -1]])

print(sigmaz)
print(np.matmul(sigmaz, sigmaz))
\end{codingbox}
	which prints \( \sigma_z = \begin{pmatrix} 1 & 0 \\ 0 & -1 \end{pmatrix} \) and \( \sigma_z^2 = \begin{pmatrix} 1 & 0 \\ 0 & 1 \end{pmatrix} \), matching the hand-calculated result.
\end{solution}
% ================================================================================ %

% ================================================================================ %
%                                    Problem 05                                    %
% ================================================================================ %
\begin{problem}
\textbf{Proof of Hermitianity}
\par Show that \( \sigma_y \) is Hermitian. Verify using Colab.
\par \( \paren{\underline{\textit{Hints:}} \text{ find } \sigma_y^{*^T} \text{ and show that it is the same as } \sigma_y} \)
\end{problem}

\begin{solution}
	By hand, the conjugate transpose of \( \sigma_y \) is obtained by transposing and then conjugating every entry:
	\[
		\sigma_y = \begin{pmatrix}
			0  & -i \\
			i & 0
		\end{pmatrix}
		\quad\Longrightarrow\quad
		\sigma_y^T = \begin{pmatrix}
			0  & i \\
			-i & 0
		\end{pmatrix}
		\quad\Longrightarrow\quad
		\sigma_y^{*^T} = \begin{pmatrix}
			0 & -i \\
			i & 0
		\end{pmatrix} = \sigma_y
	\]
	Since \( \sigma_y^{*^T} = \sigma_y \), \( \sigma_y \) is Hermitian.
	
	Verifying with Colab:
\begin{codingbox}
sigmay_dagger = sigmay.conj().T

print(sigmay_dagger)
print(np.allclose(sigmay_dagger, sigmay))
\end{codingbox}
	which prints \( \begin{pmatrix} 0 & -i \\ i & 0 \end{pmatrix} \),
	identical to \( \sigma_y \), and \texttt{True} for the equality check,
	confirming \( \sigma_y \) is Hermitian.
\end{solution}
% ================================================================================ %

% ================================================================================ %
%                                    Problem 06                                    %
% ================================================================================ %
\begin{problem}
\textbf{Operation of Paul Matrix on Bra and Ket}
\par Find \( \sigma_y \ket{0} \), \( \bra{0} \sigma_y \).
Are they the same? Why? Verify using Colab.
You can use \texttt{np.dot} for matrix-vector multiplication.
The following shows how to do \( \sigma_y \ket{0} \).
\begin{codingbox}
zerospin = np.array([[1], [0]])

print(zerospin.shape)
print(np.dot(sigmay, zerospin))
\end{codingbox}
\end{problem}

\begin{solution}
	First, \( \sigma_y \ket{0} \):
	\[
		\sigma_y \ket{0} = \begin{pmatrix}
			0 & -i \\
			i & 0
		\end{pmatrix}\begin{pmatrix}
			1 \\ 0
		\end{pmatrix} = \begin{pmatrix}
			0 \\ i
		\end{pmatrix} = i\ket{1}
	\]
	Next, \( \bra{0}\sigma_y \), where \( \bra{0} = \begin{pmatrix} 1 & 0 \end{pmatrix} \):
	\[
		\bra{0}\sigma_y = \begin{pmatrix}
			1 & 0
		\end{pmatrix}\begin{pmatrix}
			0 & -i \\
			i & 0
		\end{pmatrix} = \begin{pmatrix}
			0 & -i
		\end{pmatrix} = -i\bra{1}
	\]
	They are \textbf{not} the same: \( \sigma_y\ket{0} = i\ket{1} \) while \( \bra{0}\sigma_y = -i\bra{1} \), which differ by a sign on the imaginary coefficient.
	This is expected because, since \( \sigma_y \) is Hermitian (Problem 5), \( \bra{0}\sigma_y = \paren{\sigma_y\ket{0}}^{*^T} \) is the \emph{conjugate} transpose, not just the transpose, of \( \sigma_y\ket{0} \); as \( \sigma_y\ket{0} = i\ket{1} \) has a nonzero imaginary coefficient, conjugation flips its sign.

	Verifying with Colab:
\begin{codingbox}
zerospin = np.array([[1], [0]])
print(zerospin.shape)

sigma_y_zero = np.dot(sigmay, zerospin)
print(sigma_y_zero)

bra_zero = zerospin.conj().T
bra_zero_sigma_y = np.dot(bra_zero, sigmay)
print(bra_zero_sigma_y)
\end{codingbox}
	which prints \( \begin{pmatrix} 0 \\ i \end{pmatrix} \) for \( \sigma_y\ket{0} \) and \( \begin{pmatrix} 0 & -i \end{pmatrix} \) for \( \bra{0}\sigma_y \), matching the hand-calculated result.
\end{solution}
% ================================================================================ %

% ================================================================================ %
%                                     Appendix                                     %
% ================================================================================ %

\myappendix{code/assignment-03-code.pdf}{Assignment 03 | Code Solution}

% ================================================================================ %

\end{document}
